% STATUS: written 2026-11-06 (English gate clean; sources read in full as recorded below).
% Sources read in full: Steinitz 1910 (J. reine angew. Math. 137, 167--309; Sections 9 and 13 in the text layer; the definition of elements and extensions of the first kind in Section 13, p. 231), Teichmueller 1936 (J. reine angew. Math. 175, 89--99; introduction), Jacobson 1944 (Amer. J. Math. 66, 645--648).
\section[Derivations of fields and separability: Steinitz, Teichm\"uller, Jacobson][Fields and separability]{Derivations of fields and separability: Steinitz, Teichm\"uller, Jacobson}\label{sec:origins:fields-and-separability}

The fields of the last section have characteristic zero. In characteristic $p > 0$ a nonzero polynomial can have zero derivative. For
example, $x^p$ has zero derivative, because the polynomials with zero
derivative are exactly the polynomials in $x^p$
\cite[\S 9, p.~213]{Steinitz1910}. This fact first made
derivations of fields necessary. This section goes from the formal derivative
of Steinitz to the Galois theory of derivations of Jacobson.

The story starts in 1910. The memoir \emph{Algebraische Theorie der K\"orper}~\cite{Steinitz1910}
founded the abstract theory of fields. It is written in German. We follow it
in our translation. Section~9 is called \emph{Differenzierung}. Let
$f(x) = a_0 + a_1x + \cdots + a_nx^n$ be a polynomial over any field.
Steinitz defines its derivative \emph{rein formal} (purely formally):
\[
  f'(x) = a_1 + 2a_2x + \cdots + na_nx^{n-1}.
\]
From this formal definition Steinitz derives the sum rule and the product
rule for derivatives~\cite[\S 9, pp.~209--210]{Steinitz1910}. He then shows the following.
In characteristic $p$, we have $f'(x) = 0$ exactly when $f$ is a polynomial
in $x^p$~\cite[\S 9, p.~213]{Steinitz1910}.

Steinitz puts this fact to work at once. His Section~13 uses the derivative to separate two kinds of algebraic extension.
An irreducible polynomial is \emph{of the first kind} if it has only simple
roots~\cite[\S 13, p.~231]{Steinitz1910}. An algebraic extension is of the first kind if all its elements are.
Over a perfect field every extension is of the first kind. In particular,
this holds in characteristic $0$. Now let $g$ be irreducible over a field of
characteristic $p$. Suppose that $g$ is not a polynomial in $x^p$. Then
$g'(x) \neq 0$. So $g'$ and $g$ have no common factor. Hence $g$ has only
simple roots~\cite[\S 13, pp.~231--232]{Steinitz1910}:
\begin{originalquote}{German}{Hat zweitens der Exponent $f$ seinen kleinsten Wert $0$, so ist $g(x)$ nicht mehr ganze rationale Funktion von $x^p$, die Ableitung $g'(x)$ verschwindet nicht, kann also mit $g(x)$ keinen gemeinsamen Faktor haben, somit hat die Funktion $g(x)$ lauter einfache Nullstellen, sie ist von erster Art.}
If, secondly, the exponent $f$ has its least value $0$, then $g(x)$ is no
longer an integral rational function of $x^p$, the derivative $g'(x)$ does not
vanish, and so it can have no common factor with $g(x)$; hence the function
$g(x)$ has only simple roots, it is of the first kind.
\end{originalquote}
\noindent\cite[\S 13, p.~232]{Steinitz1910}. In present terms: an
irreducible polynomial is separable if and only if its derivative is nonzero.
The first kind of Steinitz is our separable kind.

The next step came in the 1930s. Hasse had built higher derivatives in algebraic function fields of
characteristic $p$. He used the formulas of analysis for implicit functions
(Section~\ref{sec:origins:higher-derivations}). Teichm\"uller avoids these
formulas. His paper~\cite{Teichmueller1936} is written in German, and we
follow it in our translation. He views the function field as an inseparable
extension of a well chosen subfield. From this point of view he introduces
all the higher differential quotients. He needs no power series at the
primes. He notes that this gives \translated{a theory of differentiation for
the simplest inseparable extensions of arbitrary fields of prime
characteristic on a purely algebraic basis}{eine Differentiationstheorie
f\"ur die einfachsten inseparablen Oberk\"orper beliebiger K\"orper von
Primzahlcharakteristik auf rein algebraischer Grundlage}{German}. He meant it
as a tool for his planned work on $p$-algebras~\cite[p.~89]{Teichmueller1936}.

The last step in this line is due to Jacobson, in 1944. A purely inseparable extension has no nontrivial automorphisms. So ordinary
Galois theory says nothing about it. Jacobson replaced the automorphisms by
derivations~\cite{Jacobson1944}. Let $K$ be a field of characteristic $p$.
A finite extension $K/k$ is \emph{purely inseparable of exponent one} when
$x^p \in k$ for every $x \in K$. Write $\operatorname{Der}_k(K)$ for the
derivations of $K$ that vanish on $k$.

We use the terms of Jacobson \cite[Theorem~1, p.~207]{Jacobson1937}. A set of derivations of $K$ is a
\emph{restricted $K$-Lie ring} if it has three properties. It is a vector
space over $K$. It is closed under commutators. It is closed under $p$-th
powers. It is \emph{finite} if its dimension over $K$ is finite. By
Section~\ref{sec:origins:lie-theory}, $\operatorname{Der}_k(K)$ is a
restricted $K$-Lie ring.

\begin{theorem}[{\cite[p.~645]{Jacobson1944}}]\label{thm:origins:jacobson-galois}
Let $K/k$ be a finite, purely inseparable extension of exponent one in
characteristic $p > 0$. Then the only elements of $K$ killed by every
derivation in $\operatorname{Der}_k(K)$ are the elements of $k$. Moreover,
let $\Delta$ be a finite restricted $K$-Lie ring of derivations of a field
$K$ of characteristic $p$, and let $k$ be its field of constants. Then $K/k$
is finite and purely inseparable of exponent one, and
$\Delta = \operatorname{Der}_k(K)$.
\end{theorem}

What does Theorem~\ref{thm:origins:jacobson-galois} say? Jacobson calls the result an order anti-isomorphism~\cite[p.~645]{Jacobson1944}.
It matches two kinds of objects. On one side are certain subfields $k$ of
$K$. On the other side are certain sets of derivations of $K$. A field $k$
goes to the set $\operatorname{Der}_k(K)$. A set of derivations goes back to
its field of constants. A larger field gives a smaller set of derivations.

The theorem has a history in the work of Jacobson himself. He already had a weaker result in his paper of 1937
(Section~\ref{sec:origins:lie-theory}). There he matched the fields between
$k$ and $K$ with the restricted Lie rings inside
$\operatorname{Der}_k(K)$~\cite[p.~220]{Jacobson1937}. In 1944 he can start
from either side: from the derivations, or from the base field. He compares
this step with the work of Artin and Baer in ordinary Galois
theory~\cite[p.~645]{Jacobson1944}.

So derivations now play the part of automorphisms. In ordinary Galois theory
a group of automorphisms has a fixed field. Here a set of derivations has a
field of constants.

Characteristic $p$ caused one more problem, this time for higher derivatives.
The next section tells how Hasse and Schmidt solved it.
